\documentclass[11pt,a4paper]{article}
\usepackage[T1]{fontenc}
\usepackage[utf8]{inputenc}
\usepackage{newtxtext}
\usepackage{amsmath,amsthm,mathtools}
\usepackage{newtxmath}
\usepackage[margin=28mm,headheight=15pt]{geometry}
\usepackage{microtype}
\usepackage{booktabs,tabularx,array}
\usepackage[dvipsnames]{xcolor}
\usepackage{enumitem,etoolbox}
\usepackage{fancyhdr}
\usepackage{hyperref}
\hypersetup{colorlinks=true,linkcolor=MidnightBlue,citecolor=MidnightBlue,
 urlcolor=MidnightBlue,pdftitle={Classical Complex Theorems over the Quaternions},
 pdfauthor={OpenAI},pdfsubject={Slice regularity, Cauchy--Fueter regularity, and proofs of classical analogues}}
\urlstyle{same}
\usepackage{bookmark}
\numberwithin{equation}{section}
\newtheorem{theorem}{Theorem}[section]
\newtheorem{proposition}[theorem]{Proposition}
\newtheorem{lemma}[theorem]{Lemma}
\newtheorem{corollary}[theorem]{Corollary}
\theoremstyle{definition}
\newtheorem{definition}[theorem]{Definition}
\newtheorem{example}[theorem]{Example}
\theoremstyle{remark}
\newtheorem{remark}[theorem]{Remark}
\newcommand{\R}{\mathbb R}
\newcommand{\C}{\mathbb C}
\newcommand{\HH}{\mathbb H}
\newcommand{\SSS}{\mathbb S}
\newcommand{\BB}{\mathbb B}
\newcommand{\ii}{\mathbf i}
\newcommand{\jj}{\mathbf j}
\newcommand{\kk}{\mathbf k}
\newcommand{\Id}{\operatorname{id}}
\newcommand{\Rea}{\operatorname{Re}}
\newcommand{\Ima}{\operatorname{Im}}
\newcommand{\SR}{\mathcal{SR}}
\newcommand{\cD}{\mathcal D}
\newcommand{\lap}{\Delta_4}
\newcommand{\Norm}{\mathcal N}
\newcommand{\cI}{\mathcal I}
\newcommand{\cC}{\mathcal C}
\newcommand{\sph}{\partial_{\mathrm{sph}}}
\newcommand{\pc}{\partial_c}
\newcommand{\Res}{\operatorname{Res}}
\newcommand{\ord}{\operatorname{ord}}
\newcommand{\dd}{\,\mathrm d}
\newcommand{\qbar}[1]{\overline{#1}}
\newcommand{\abs}[1]{\lvert #1\rvert}
\newcommand{\norm}[1]{\lVert #1\rVert}
\DeclareMathOperator{\supp}{supp}
\setlist{itemsep=2pt,topsep=4pt}
\setlength{\parindent}{1.2em}
\setlength{\parskip}{2pt}
\setlength{\emergencystretch}{2em}
\pagestyle{fancy}
\fancyhf{}
\fancyhead[L]{\small Classical complex theorems over the quaternions}
\fancyhead[R]{\small\thepage}
\renewcommand{\headrulewidth}{0.3pt}
\renewcommand{\sectionmark}[1]{\markboth{#1}{}}
\title{\vspace{-12mm}\textbf{Classical Complex Theorems\\over the Quaternions}\\[5pt]
\large Slice-regular and Cauchy--Fueter analogues, with proofs}
\author{An expository mathematical paper}
\date{20 September 2026}
\begin{document}
\maketitle
\thispagestyle{empty}
\begin{abstract}
Complex differentiability combines properties that separate when the scalars are
quaternions. We develop two complementary replacements and prove a substantial
collection of classical analogues. For slice-regular functions, we establish
representation and Cauchy formulas, Taylor and Laurent expansions, identity and
maximum principles, Liouville and Schwarz theorems, a regular Schwarz--Pick
inequality, a fundamental theorem of algebra, a multiplicity-sensitive argument
principle and Rouch\'e theorem, minimum-modulus and open-mapping results, Montel
and Hurwitz theorems, a Casorati--Weierstrass theorem, and an intrinsic Riemann
mapping theorem. For Cauchy--Fueter-regular functions, we derive the
four-dimensional Cauchy--Pompeiu formula directly from the divergence theorem,
then prove derivative estimates, mean-value and maximum principles, Liouville,
Morera, removable-singularity, and residue theorems. A direct proof of Fueter's
mapping theorem explains how the Laplacian connects the theories; an explicit
local inverse is obtained by integrating two closed one-forms. Throughout,
multiplication order, the role of the real axis, spherical zeros, and the limits
of the complex analogy are made explicit.
\end{abstract}
\vspace{3pt}
\noindent\textbf{Keywords.} Quaternionic analysis; slice regularity; Cauchy--Fueter
operator; spherical zeros; regular products; Fueter mapping theorem.

\medskip
\noindent\textbf{Scope and prerequisites.} This is a synthesis of established
mathematics, not a claim of new priority. All quaternionic results stated as
theorems, propositions, or lemmas are proved. We use the usual scalar theorems of
one-variable complex analysis, including Cauchy's theorem, the maximum principle,
Laurent's theorem, Montel, Hurwitz, and the normalized Riemann mapping theorem.
We also use elementary real analysis, the divergence theorem, and the existence
of smooth mollifiers. Neither quaternionic function theory nor Clifford
analysis is assumed.
\newpage
\begingroup
\small
\setlength{\parskip}{0pt}
\makeatletter
\patchcmd{\l@section}{1.0em}{0.45em}{}{}
\makeatother
\tableofcontents
\endgroup
\newpage

\section{Why there are two theories}
\label{sec:intro}
For complex-valued functions of one complex variable, differentiability, local
power-series expansions, the Cauchy--Riemann equation, and several geometric
properties fit into a single theory. Over the noncommutative division algebra
$\HH$, the most literal translations do not fit together. Slice regularity
preserves power series in one quaternionic variable. Cauchy--Fueter regularity
preserves a first-order elliptic differential equation and a boundary integral
in four real dimensions. The theories originate in different constructions;
modern foundations are given in
\cite{GS,GP,CGS,Sudbery}. The relation between them is a differential transform,
not an identification of their function spaces.

\subsection{Algebra and conventions}
Write
\[
 q=x_0+\ii x_1+\jj x_2+\kk x_3,
 \qquad \ii^2=\jj^2=\kk^2=\ii\jj\kk=-1.
\]
Quaternionic conjugation reverses products:
$\overline{ab}=\bar b\bar a$. Set
\[
 \Rea q=\frac{q+\bar q}{2},\qquad
 |q|^2=q\bar q=\bar q q,\qquad q^{-1}=\frac{\bar q}{|q|^2}\quad(q\ne0).
\]
In particular, $|ab|=|a||b|$. The sphere of imaginary units and its associated
complex planes are
\[
 \SSS=\{I\in\HH:I^2=-1\},\qquad \C_I=\R+I\R.
\]
Every nonreal quaternion has a unique expression $x+Iy$ with $y>0$ and
$I\in\SSS$. Its conjugacy class is the two-sphere
\[
 [q]=x+y\SSS;
\]
for real $q$, $[q]=\{q\}$. Conjugation by a nonzero quaternion preserves the real
part and norm and acts on these spheres. We use $\iota$ for the imaginary unit of
an auxiliary, \emph{commutative} complex field; $\iota$ is not a chosen element of
$\SSS$.

All slice-regular functions below are \emph{left} slice regular. Their power
series have coefficients on the right:
\[
 f(q)=\sum_{n\ge0}q^n a_n.
\]
Constants may always be multiplied on the right. A constant on the left requires
care. All integration of quaternion-valued functions means integration of four
real components, with the displayed order of quaternionic factors retained.

\subsection{The literal derivative is too restrictive}
\begin{proposition}[Rigidity of a direction-independent quotient]
\label{prop:rigidity}
Let $U\subseteq\HH$ be connected and open, and let $f\in C^2(U,\HH)$.
If, for every $q\in U$, the limit
\[
 a(q)=\lim_{h\to0}h^{-1}\bigl(f(q+h)-f(q)\bigr)
\]
exists independently of the direction of $h$, then $f(q)=q a+b$ for fixed
$a,b\in\HH$.
\end{proposition}
\begin{proof}
Put $e_0=1$, $e_1=\ii$, $e_2=\jj$, and $e_3=\kk$.
Taking increments $h=t e_\mu$ gives
$\partial_\mu f=e_\mu a$ and $a=\partial_0 f$.
Equality of mixed derivatives with $\partial_0$ implies
$\partial_\mu a=e_\mu\partial_0 a$. Equality of the derivatives with indices
$1,2$ now gives
\[
 \ii\jj\,\partial_0a=\jj\ii\,\partial_0a.
\]
Since $\ii\jj-\jj\ii=2\kk$ is invertible, $\partial_0a=0$, and therefore all
first derivatives of $a$ vanish. Connectedness makes $a$ constant, and all
first derivatives of $f(q)-qa$ vanish as well.
\end{proof}
The analogous quotient with $h^{-1}$ on the right produces $f(q)=aq+b$.
The $C^2$ assumption in the proposition is explicit; no lower-regularity
rigidity theorem is needed here. In particular, this definition excludes most
polynomials. We therefore choose a different structure rather than trying to
force the complex difference quotient to do everything.

\section{Slice regularity, representation, and differential structure}
\label{sec:slices}
\subsection{Domains and the splitting lemma}
Let $D\subseteq\C$ be a connected open set invariant under complex conjugation
and meeting $\R$. Its \emph{circularization} is
\[
 \Omega_D=\{x+Iy:x+\iota y\in D,\ I\in\SSS\}.
\]
It is an axially symmetric slice domain: it contains a whole conjugacy sphere
whenever it contains one point of that sphere, and every section
$\Omega_D\cap\C_I$ is connected. Unless another domain is explicitly specified,
this is the setting for slice-regular theorems.

\begin{definition}
A $C^1$ function $f:\Omega_D\to\HH$ is \emph{slice regular} if, for each
$I\in\SSS$, its restriction $f_I$ satisfies
\begin{equation}
 \bar\partial_I f_I
 =\frac12(\partial_x+I\partial_y)f_I=0.
 \label{eq:CRslice}
\end{equation}
Its slice, or Cullen, derivative is
$\pc f=\frac12(\partial_x-I\partial_y)f=\partial_x f$.
\end{definition}

\begin{lemma}[Splitting]
\label{lem:splitting}
Fix orthogonal $I,J\in\SSS$. Every quaternion-valued function on $\C_I$ can
be written uniquely as $F+GJ$, where $F,G$ take values in $\C_I$.
It satisfies \eqref{eq:CRslice} if and only if $F,G$ are complex holomorphic.
Moreover, $|F+GJ|^2=|F|^2+|G|^2$.
\end{lemma}
\begin{proof}
The real basis $1,I,J,IJ$ gives the decomposition. Substituting it into
\eqref{eq:CRslice}, the two $\C_I$ components vanish independently and give
ordinary Cauchy--Riemann equations. The norm identity follows from the
orthogonality of the four basis elements.
\end{proof}

\begin{theorem}[Identity principle]
\label{thm:identity}
If two slice-regular functions on $\Omega_D$ agree on a subset of one
$\C_I$-section having an accumulation point in that section, they agree on
$\Omega_D$.
\end{theorem}
\begin{proof}
Apply scalar uniqueness to both components of their difference in the splitting
lemma. The difference vanishes on the entire $I$-section and hence on a real
interval contained in $D$. Its restriction to any other section has those same
real zeros. Applying the splitting lemma and scalar uniqueness again proves
the claim on every section.
\end{proof}
Notice the phrase ``one section.'' An arbitrary accumulation of zeros in
$\HH$ does not suffice; a whole two-sphere may be a zero set.

\subsection{Representation and stem functions}
\begin{theorem}[Representation formula]
\label{thm:representation}
For every slice-regular $f$ on $\Omega_D$ and every $I,K\in\SSS$,
\begin{equation}
 f(x+Ky)=\frac{1-KI}{2}f(x+Iy)+\frac{1+KI}{2}f(x-Iy).
 \label{eq:representation}
\end{equation}
Equivalently, there are uniquely determined quaternion-valued functions
$A(x,y),B(x,y)$ such that
\begin{equation}
 f(x+Ky)=A(x,y)+KB(x,y),
 \label{eq:AB}
\end{equation}
where $A$ is even and $B$ is odd in $y$, and
\begin{equation}
 A_x=B_y,\qquad A_y=-B_x.
 \label{eq:stemCR}
\end{equation}
\end{theorem}
\begin{proof}
On one fixed section put $h_\pm(x,y)=f(x\pm Iy)$ and define
\[
 A=\frac{h_++h_-}{2},\qquad B=-\frac I2(h_+-h_-).
\]
The parity statements are immediate. Equation \eqref{eq:CRslice} gives
$(h_+)_y=I(h_+)_x$ and $(h_-)_y=-I(h_-)_x$, which imply
\eqref{eq:stemCR}. Thus $A+KB$ is holomorphic on each $K$-section.
At $y=0$ it equals $f(x)$, so scalar uniqueness on that section makes it equal
to $f(x+Ky)$ everywhere. Expanding $A+KB$ gives
\eqref{eq:representation}. Evaluating at $K$ and $-K$ proves uniqueness.
\end{proof}

Introduce the complexified algebra
\[
 \HH_{\C}=\{a+\iota b:a,b\in\HH\},\qquad
 (a+\iota b)(c+\iota d)=ac-bd+\iota(ad+bc).
\]
Here $\iota$ commutes with all quaternionic elements. Let
$\kappa(a+\iota b)=a-\iota b$; this conjugation is different from quaternionic
conjugation. The function
\[
 \mathcal F(x+\iota y)=A(x,y)+\iota B(x,y)
\]
is holomorphic with values in the four-dimensional complex vector space
$\HH_{\C}$ and satisfies
\begin{equation}
 \mathcal F(\bar z)=\kappa(\mathcal F(z)).
 \label{eq:stem-symmetry}
\end{equation}
Such a function is called a \emph{holomorphic stem}. Conversely, a holomorphic
stem induces a slice-regular function by \eqref{eq:AB}; denote this operation by
$f=\cI(\mathcal F)$. Away from the real axis this follows directly from
\eqref{eq:stemCR}. Near a real point the complex Taylor coefficients of
$\mathcal F$ belong to $\HH$, by \eqref{eq:stem-symmetry}, and its induced
convergent power series is real analytic. This also justifies the $C^1$
regularity at real points in the converse construction.

A slice function is \emph{intrinsic}, or slice preserving, when its $A,B$ are
real valued. Its values on each $\C_I$ lie in that same plane. In particular,
real-coefficient power series are intrinsic. The stem formalism, developed in
\cite{GP}, makes the distinction between an auxiliary complex scalar and a
quaternionic imaginary unit precise.

\subsection{Two derivatives at a nonreal point}
\begin{proposition}[Real differential]
\label{prop:differential}
Let $q=x+Iy$, $y>0$, and decompose $v\in\HH$ orthogonally as
$v=v_\parallel+v_\perp$ with $v_\parallel\in\C_I$ and
$v_\perp\in\C_I^\perp$. Then
\begin{equation}
 df_q(v)=v_\parallel\pc f(q)+v_\perp\sph f(q),\qquad
 \sph f(q)=\frac{B(x,y)}{y}
 =(2Iy)^{-1}\bigl(f(q)-f(\bar q)\bigr).
 \label{eq:differential}
\end{equation}
At a real point $x$, $df_x(v)=v\pc f(x)$.
\end{proposition}
\begin{proof}
Within $\C_I$, the Cauchy--Riemann equation gives
$\partial_yf=I\partial_xf$, which proves the parallel formula. In a
perpendicular direction, the real part and the length of the imaginary part
have zero first variation, whereas $I$ varies by $v_\perp/y$. Differentiate
\eqref{eq:AB}. At a real point use the real-centered power series described
above; its linear part is $v\pc f(x)$.
\end{proof}
For $f(q)=q^2$, these two right multipliers are $2q$ and $2x$.
They need not have equal norms, so slice regularity does not imply
four-dimensional conformality.

\section{Cauchy formulas, Taylor series, and Morera}
\label{sec:cauchy-slice}
\subsection{The kernel and its multiplication order}
For $q\notin[s]$, set
\begin{equation}
 \cC(q,s)=S_L^{-1}(s,q)
 =-\bigl(q^2-2\Rea(s)q+|s|^2\bigr)^{-1}(q-\bar s).
 \label{eq:kernel}
\end{equation}
The inverse on the right is the ordinary quaternionic inverse of the displayed
quadratic value. For fixed $s$, this is a slice-regular function of $q$ away
from $[s]$: it is induced by the holomorphic stem obtained by replacing $q$
with the central complex variable $z$ in \eqref{eq:kernel}. When $q,s\in\C_I$,
commutativity gives $\cC(q,s)=(s-q)^{-1}$.

For $|q|<|s|$ the kernel also has the expansion
\begin{equation}
 \cC(q,s)=\sum_{n\ge0}q^n s^{-n-1}.
 \label{eq:kernelseries}
\end{equation}
Indeed, the absolutely convergent sum $R$ satisfies $Rs-qR=1$.
Using $s^2-2\Rea(s)s+|s|^2=0$ then gives
\[
 0=R\bigl(s^2-2\Rea(s)s+|s|^2\bigr)
   =\bigl(q^2-2\Rea(s)q+|s|^2\bigr)R+q-\bar s,
\]
which proves \eqref{eq:kernelseries}. This is not a license to commute $q$
with $s$.

\begin{theorem}[Slice Cauchy formula]
\label{thm:cauchy-slice}
Suppose $D$ is bounded with piecewise $C^1$ boundary and $f$ is slice regular
on an axially symmetric slice domain containing $\overline{\Omega_D}$.
Give $\partial D_I$ its positive complex orientation, including the usual
negative orientation on holes, and set $\mathrm ds_I=-I\,\mathrm ds$.
Then, for all $q\in\Omega_D$,
\begin{equation}
 f(q)=\frac1{2\pi}\int_{\partial D_I}
          \cC(q,s)\,\mathrm ds_I\,f(s).
 \label{eq:cauchy-slice}
\end{equation}
The answer is independent of $I$.
\end{theorem}
\begin{proof}
On the $I$-section, use the splitting lemma and apply the ordinary Cauchy
formula to both complex components. The normalization $\mathrm ds_I=-I\,
\mathrm ds$ puts the complex factor $1/I$ in the correct position.
For arbitrary $q$, both sides of \eqref{eq:cauchy-slice} are slice functions
of $q$ and obey \eqref{eq:representation}. This is true for the integral
because it holds for every kernel and the factors $\mathrm ds_I f(s)$ are
on the right. Equality on the $I$-section therefore proves equality everywhere.
The kernel has no singularities in $\Omega_D$ when $s\in\partial D_I$,
because $D$ is conjugation invariant. Independence of $I$ follows from the
same equality with $f(q)$.
\end{proof}
This kernel and formula are standard results of \cite{CGS}; the proof here
makes explicit how the representation formula replaces commutativity.

\begin{corollary}[Cauchy's theorem]
Under the corresponding holomorphy assumptions,
\[
 \int_{\partial U_I}\mathrm ds_I\, f(s)=0
\]
for any relatively compact planar region $U_I$ in a section on which $f$
is regular up to the boundary.
\end{corollary}
\begin{proof}
Apply scalar Cauchy's theorem to the two components in the splitting lemma.
\end{proof}

\subsection{Taylor expansion and estimates}
\begin{theorem}[Taylor theorem at real centers]
\label{thm:taylor}
Let $c\in\R$, and let $f$ be slice regular on $B(c,R)$. There are unique
$a_n\in\HH$ such that
\begin{equation}
 f(q)=\sum_{n=0}^{\infty}(q-c)^n a_n,\qquad
 a_n=\frac{\pc^n f(c)}{n!}.
 \label{eq:taylor}
\end{equation}
The series converges absolutely and uniformly on smaller closed balls.
For $0<r<R$,
\begin{align}
 a_n&=\frac1{2\pi}\int_{|s-c|=r,\ s\in\C_I}
          (s-c)^{-n-1}\,\mathrm ds_I\,f(s),\label{eq:coefficient}\\
 |a_n|&\le \frac{M_I(r)}{r^n},\qquad
 M_I(r)=\max_{|s-c|=r,\ s\in\C_I}|f(s)|.\label{eq:cauchy-estimate}
\end{align}
\end{theorem}
\begin{proof}
Take the usual Taylor expansions of both components on one section. Their
quaternionic coefficients combine into $a_n$, and real differentiation at
$c$ identifies them independently of the section. Scalar Cauchy's formula
gives \eqref{eq:coefficient}. Its norm estimate is
\[
 \frac1{2\pi}(2\pi r)r^{-n-1}M_I(r)=r^{-n}M_I(r).
\]
For $|q-c|\le\rho<r$, the resulting power series is bounded termwise by a
convergent geometric series. It defines a slice-regular function, agrees with
$f$ on the chosen section, and hence agrees everywhere by
Theorem~\ref{thm:identity}. Uniqueness follows by real differentiation.
\end{proof}
Consequently every slice-regular function is real analytic: the stem description
proves this off the real axis, and \eqref{eq:taylor} proves it at the real axis.
Ordinary powers centered at a nonreal quaternion are not asserted to give a
Taylor theorem on an ordinary Euclidean ball.

\begin{theorem}[Morera]
\label{thm:morera-slice}
Let $f:\Omega_D\to\HH$ be continuous. Suppose that for every $I\in\SSS$ and
every triangle $T$ with closure in the $I$-section,
\[
 \int_{\partial T}\mathrm ds_I\,f(s)=0.
\]
Then $f$ is slice regular.
\end{theorem}
\begin{proof}
Split $f_I$ into its two continuous complex components. Scalar Morera makes
each component holomorphic. Construct $A,B$ from one section as in the proof
of Theorem~\ref{thm:representation}. They give a real-analytic slice-regular
extension. On every other section that extension and $f$ are holomorphic and
agree on a real interval, so they coincide.
\end{proof}

\section{Maximum principles, growth, and Schwarz's lemma}
\label{sec:maximum}
\begin{theorem}[Maximum-modulus principle]
\label{thm:max}
If $f\in\SR(\Omega_D)$ and $|f|$ has a local maximum at an interior point,
then $f$ is constant.
\end{theorem}
\begin{proof}
Let the maximum occur at $p\in\C_I$, and put $M=|f(p)|$. If $M=0$, the
function vanishes on a neighborhood and the identity principle applies.
Otherwise, write $u=f(p)/M$ and $g=f\bar u$. Then $g$ is regular, $g(p)=M$,
and $|g|=|f|$. Split $g_I=F+GJ$. Near $p$, $|F|\le|g|\le M$ and
$F(p)=M$. Scalar maximum modulus gives $F\equiv M$ locally. The identity
$|g|^2=|F|^2+|G|^2\le M^2$ forces $G=0$ there. The identity principle
makes $f$ constant on $\Omega_D$.
\end{proof}
In particular, for a function regular on $\Omega_D$, the maximum of its
norm on any closed Euclidean ball contained in $\Omega_D$ is attained on
the boundary.

\begin{theorem}[Liouville and polynomial growth]
\label{thm:liouville-slice}
A bounded entire slice-regular function is constant. More generally, if
\[
 |f(q)|\le C(1+|q|)^\alpha\qquad(q\in\HH),\quad \alpha\ge0,
\]
then $f$ is a right-coefficient polynomial of degree at most $\lfloor\alpha\rfloor$.
\end{theorem}
\begin{proof}
Expand at zero. For every $r>0$, \eqref{eq:cauchy-estimate} yields
$|a_n|\le C(1+r)^\alpha/r^n$. Let $r\to\infty$ for each integer $n>\alpha$.
The corresponding coefficient vanishes. Boundedness is the case $\alpha=0$.
\end{proof}

\begin{theorem}[Higher-order Schwarz lemma]
\label{thm:schwarz}
Let $f:\BB\to\BB$ be slice regular, where $\BB=B(0,1)$, and suppose the
first $m$ coefficients in its Taylor series vanish, with $m\ge1$.
Then
\begin{equation}
 |f(q)|\le |q|^m,\qquad \frac{|\pc^m f(0)|}{m!}\le1.
 \label{eq:schwarz}
\end{equation}
Equality in the first inequality at a nonzero point, or in the second
inequality, occurs exactly when $f(q)=q^m u$ for a fixed $|u|=1$.
\end{theorem}
\begin{proof}
Write $f(q)=q^m g(q)$ using its power series; $g$ is regular on $\BB$.
For $|q|=r<1$, $|g(q)|\le r^{-m}$. Maximum modulus on $B(0,r)$ and then
$r\uparrow1$ give $|g(q)|\le1$ everywhere. This proves the two estimates.
Either equality hypothesis forces $|g|$ to attain the value $1$ at an interior
point. Theorem~\ref{thm:max} makes $g$ a constant of unit norm. The converse
is immediate from multiplicativity of the quaternionic norm.
\end{proof}
These basic analogues belong to the original slice-regular theory
\cite{GS}. Their proofs illustrate a useful principle: use scalar holomorphic
components for linear analytic questions, but preserve the quaternionic norm
when a sharp scalar bound is required.

\section{Regular multiplication and the geometry of zeros}
\label{sec:algebra}
Pointwise multiplication does not preserve slice regularity. For example, the
constant function $\ii$ and the identity function are regular, but on the
$\jj$-section the product $\ii q=\ii x+\kk y$ satisfies
\[
 (\partial_x+\jj\partial_y)(\ii q)=2\ii\ne0.
\]
The appropriate replacement is multiplication before evaluating the stem.

\subsection{The regular product}
\begin{definition}
If $f=\cI(\mathcal F)$ and $g=\cI(\mathcal G)$, define
\[
 f*g=\cI(\mathcal F\mathcal G).
\]
Thus, when $f(x+Iy)=A+IB$ and $g(x+Iy)=C+ID$,
\begin{equation}
 (f*g)(x+Iy)=AC-BD+I(AD+BC).
 \label{eq:starAB}
\end{equation}
\end{definition}
Stem symmetry is preserved because $\kappa$ is an algebra automorphism;
holomorphy is preserved because multiplication is complex bilinear. Therefore
$*$ is an associative product on slice-regular functions. For power series,
\begin{equation}
 \left(\sum_{n\ge0}q^n a_n\right)*
 \left(\sum_{n\ge0}q^n b_n\right)
 =\sum_{n\ge0}q^n\sum_{k=0}^n a_k b_{n-k}.
 \label{eq:convolution}
\end{equation}
This follows by multiplying absolutely convergent stem series on compact
subsets. Intrinsic functions have central stems: if $h$ is intrinsic, then
$h*f=f*h$, and this regular product equals the \emph{pointwise product
$h(q)f(q)$}. It need not equal $f(q)h(q)$.

\begin{proposition}[Evaluation of a regular product]
\label{prop:evaluation}
For any $q\in\Omega_D$,
\begin{equation}
 (f*g)(q)=
 \begin{cases}
 f(q)\,g\bigl(f(q)^{-1}qf(q)\bigr),&f(q)\ne0,\\
 0,&f(q)=0.
 \end{cases}
 \label{eq:evaluation}
\end{equation}
\end{proposition}
\begin{proof}
Write $q=x+Iy$ and $u=f(q)=A+IB$. For $u\ne0$, set $I'=u^{-1}Iu$.
The proposed right side is
\[
 u(C+I'D)=uC+IuD=AC-BD+I(BC+AD),
\]
which is \eqref{eq:starAB}. The same expression is $uC+IuD=0$ when $u=0$.
\end{proof}
The argument of the second factor can rotate within $[q]$. This is precisely
why roots of successive factors need not be roots of their product.

\subsection{Conjugation and symmetrization}
Extend quaternionic conjugation complex linearly to the stem algebra:
$(a+\iota b)^c=\bar a+\iota\bar b$. It is an antiautomorphism. Define
\[
 f^c=\cI(\mathcal F^c),\qquad f^s=f*f^c=f^c*f.
\]
The equality of the two products follows from the following direct calculation:
\begin{align}
 \Norm_f(z):=\mathcal F(z)\mathcal F(z)^c
 &=|A|^2-|B|^2+2\iota\Rea(A\bar B),\label{eq:normstem}\\
 f^s(x+Iy)&=|A|^2-|B|^2+2I\Rea(A\bar B).
 \label{eq:symmetrization}
\end{align}
The stem $\Norm_f$ is complex scalar valued and holomorphic; hence $f^s$ is
intrinsic. In a right-coefficient power series,
\[
 f^c(q)=\sum q^n\bar a_n,\qquad
 f^s(q)=\sum q^n\sum_{k=0}^n a_k\bar a_{n-k},
\]
and each coefficient of $f^s$ is real. In particular,
\begin{equation}
 \Norm_f(x)=|f(x)|^2\quad(x\in D\cap\R).
 \label{eq:normonreal}
\end{equation}
For general $q$, $f^s(q)$ is \emph{not} $|f(q)|^2$.
For instance, $f(q)=q-\ii$ has $f^s(q)=q^2+1$; at $q=\jj$ these two
quantities are $0$ and $2$, respectively.

\begin{proposition}[Multiplicativity of the scalar norm stem]
\label{prop:norm-multiplicative}
For slice-regular $f,g$,
\begin{equation}
 \Norm_{f*g}=\Norm_f\Norm_g,
 \qquad (f*g)^s=f^s g^s.
 \label{eq:norm-multiplicative}
\end{equation}
If $f\not\equiv0$, then $\Norm_f\not\equiv0$.
\end{proposition}
\begin{proof}
Since $\Norm_g$ is central,
\[
 (\mathcal F\mathcal G)(\mathcal F\mathcal G)^c
 =\mathcal F\mathcal G\mathcal G^c\mathcal F^c
 =\Norm_g\mathcal F\mathcal F^c=\Norm_g\Norm_f.
\]
If $\Norm_f$ vanished identically, \eqref{eq:normonreal} would make $f$ vanish
on a real interval, and the identity principle would give $f\equiv0$.
\end{proof}

\subsection{Isolated zeros and spherical zeros}
\begin{theorem}[Zero geometry and detection by symmetrization]
\label{thm:zeros}
On a nonreal sphere $S=x+y\SSS$, a slice function has either no zero,
exactly one zero, or is zero everywhere. Moreover,
\begin{equation}
 f^s\big|_S=0
 \quad\Longleftrightarrow\quad
 f\text{ has a zero on }S.
 \label{eq:zero-detection}
\end{equation}
For real $x$, $f^s(x)=0$ if and only if $f(x)=0$.
If $f$ is regular and not identically zero, its zero-bearing conjugacy classes
are locally finite in $\Omega_D$. Thus its zeros consist of isolated real or
nonreal points and isolated two-spheres.
\end{theorem}
\begin{proof}
On $S$ write $f(x+Iy)=A+IB$. If $B=0$, the assertion is immediate. If
$B\ne0$, the only possible zero is at
\[
 I=-AB^{-1}.
\]
This quaternion is an imaginary unit exactly when
$|A|=|B|$ and $\Rea(A\bar B)=0$: its norm is $|A|/|B|$ and its real part
is $-\Rea(A\bar B)/|B|^2$. These are exactly the two real conditions for
\eqref{eq:symmetrization} to vanish. The real case is
\eqref{eq:normonreal}.

For a nonzero regular function, $\Norm_f$ is a nonzero scalar holomorphic
function. Its zeros are discrete in $D$. Each conjugate pair corresponds to
one nonreal conjugacy sphere, and each real zero corresponds to one real
point. A compact subset of $\Omega_D$ projects to a compact subset of $D$
under $q\mapsto\Rea q\pm\iota|\Ima q|$, so only finitely many such classes
can meet it.
\end{proof}
This description is one of the structural differences from complex analysis;
see \cite{GStZeros,GP}. For example, $q^2+1$ vanishes on all of $\SSS$.

\begin{proposition}[Division by a spherical factor]
\label{prop:spherical-division}
Let $p=x+Iy$ be nonreal and set
\[
 \Delta_p(q)=q^2-2xq+x^2+y^2.
\]
A regular function $f$ vanishes identically on $[p]$ if and only if
$f=\Delta_p*g=\Delta_p\,g$ for a regular $g$ on $\Omega_D$.
\end{proposition}
\begin{proof}
Vanishing on the sphere means $A(x,y)=B(x,y)=0$, so the stem $\mathcal F$
vanishes at $z_0=x+\iota y$ and $\bar z_0$. Divide each of its four complex
components by $(z-z_0)(z-\bar z_0)$. Both apparent singularities are removable,
and the quotient has the required stem symmetry because this denominator has
real coefficients. The converse follows because $\Delta_p$ vanishes on $[p]$.
\end{proof}

\section{Factorization, the argument principle, and Rouch\'e}
\label{sec:counting}
\subsection{Division at one point and the fundamental theorem of algebra}
\begin{lemma}[Regular division on a ball]
\label{lem:division}
Let $f(q)=\sum_{n\ge0}q^n a_n$ converge on $B(0,R)$ and let $p\in B(0,R)$.
There is a unique regular $g$ on that ball such that
\begin{equation}
 f=(q-p)*g+f(p).
 \label{eq:division}
\end{equation}
Its coefficients are
\begin{equation}
 g(q)=\sum_{n\ge0}q^n b_n,
 \qquad b_n=\sum_{k=n+1}^{\infty}p^{k-n-1}a_k.
 \label{eq:divisioncoeff}
\end{equation}
For a polynomial of degree $n\ge1$, $g$ has degree $n-1$.
\end{lemma}
\begin{proof}
Choose $|p|<r<R$. Cauchy's estimate gives $|a_k|\le M(r)r^{-k}$, so
\[
 |b_n|\le\frac{M(r)r^{-n-1}}{1-|p|/r}.
\]
Thus the series for $g$ converges on $|q|<r$, and, since $r$ can approach $R$,
on the whole ball. Directly from \eqref{eq:divisioncoeff},
$b_{n-1}-p b_n=a_n$ for $n\ge1$, while
$-p b_0=a_0-f(p)$. These are exactly the coefficients in
\eqref{eq:division}. If two quotients existed, their difference $h$ would
satisfy $(q-p)*h=0$. At every real $x\ne p$, evaluation gives $(x-p)h(x)=0$,
so the identity principle gives $h=0$. Polynomial division is the finite
version of the same computation.
\end{proof}

\begin{theorem}[Quaternionic fundamental theorem of algebra]
\label{thm:FTA}
Every nonconstant right-coefficient quaternionic polynomial has a zero in
$\HH$. If
$P(q)=q^n a_n+\cdots+a_0$, $a_n\ne0$, then
\begin{equation}
 P=(q-p_1)*\cdots*(q-p_n)\,a_n
 \label{eq:factorization}
\end{equation}
for suitable $p_1,\ldots,p_n\in\HH$.
\end{theorem}
\begin{proof}
The real-coefficient polynomial $P^s$ has degree $2n$ and leading coefficient
$|a_n|^2$. By the complex fundamental theorem of algebra its scalar stem has
a root $z$. A real root gives $P(z)=0$; a nonreal root gives a zero of $P$ on
the corresponding sphere by Theorem~\ref{thm:zeros}. Divide out one linear
regular factor by Lemma~\ref{lem:division}, and repeat on the quotient until
a constant remains. Leading coefficients identify that constant as $a_n$.
\end{proof}
The result concerns the central-variable, right-coefficient polynomial algebra.
It is not a claim about arbitrary expressions with $q$ interspersed among
quaternionic coefficients.

\begin{example}[Factors do not simply list roots]
\label{ex:camshaft}
The polynomial
\[
 P=(q-\ii)*(q-\jj)=q^2-q(\ii+\jj)+\kk
\]
satisfies $P(\ii)=0$ but $P(\jj)=2\kk\ne0$. Its symmetrization is
$(q^2+1)^2$, and it has just one pointwise zero, $\ii$: on $\SSS$ its
coefficient $B=-(\ii+\jj)$ is nonzero, so Theorem~\ref{thm:zeros} allows
only one zero; no other sphere can contain a zero. In contrast,
$(q-\ii)*(q+\ii)=q^2+1$ vanishes on an entire sphere, although it has the
same symmetrization $(q^2+1)^2$.
\end{example}

\subsection{A precise total multiplicity}
For a nonzero regular $f$, define the \emph{total multiplicity of a
zero-bearing class} by
\begin{equation}
 m_f([p])=
 \begin{cases}
 \ord_{x+\iota y}\Norm_f,&p=x+Iy,\ y>0,\\[2pt]
 \tfrac12\ord_x\Norm_f,&p=x\in\R.
 \end{cases}
 \label{eq:multiplicity}
\end{equation}
The two nonreal conjugate zeros have the same order. The real order is even:
if $f(q)=(q-x)^m u(q)$ with $u(x)\ne0$, as supplied by the Taylor theorem,
then
\[
 \Norm_f(z)=(z-x)^{2m}\Norm_u(z),\qquad \Norm_u(x)=|u(x)|^2>0.
\]
Thus \eqref{eq:multiplicity} is a positive integer. A simple isolated linear
factor has total multiplicity one; the full spherical factor $\Delta_p$ has
total multiplicity two. This convention counts algebraic mass in a conjugacy
class, not the cardinality of the pointwise zero set.

\begin{corollary}[Degree equals total multiplicity]
\label{cor:degreecount}
For a degree-$n$ polynomial,
\[
 \sum_{[p]:\ P\text{ has a zero on }[p]}m_P([p])=n.
\]
\end{corollary}
\begin{proof}
By \eqref{eq:norm-multiplicative} and \eqref{eq:factorization},
\[
 P^s(q)=|a_n|^2\prod_{j=1}^n\Delta_{p_j}(q).
\]
Each nonreal factor contributes one to the multiplicity of its conjugacy
class; each real factor is $(q-p_j)^2$ and contributes one under the
real convention in \eqref{eq:multiplicity}. Equivalently, take half the total
complex zero count of the degree-$2n$ polynomial $\Norm_P$.
\end{proof}

\subsection{Counting on circular regions}
Let $U\Subset D$ be a bounded conjugation-invariant region with finitely many
piecewise $C^1$ boundary curves. It need not be connected. Write
\[
 \Gamma_U=\{x+Iy:x+\iota y\in\partial U,\ I\in\SSS\}.
\]
Suppose $f$ has no zero on $\Gamma_U$, and let $\nu_f(U)$ be the sum of
\eqref{eq:multiplicity} over the zero-bearing classes in $\Omega_U$.

\begin{theorem}[Argument principle]
\label{thm:argument}
With the preceding assumptions,
\begin{equation}
 \nu_f(U)=\frac1{4\pi\iota}
    \int_{\partial U}\frac{\Norm_f'(z)}{\Norm_f(z)}\dd z.
 \label{eq:argument}
\end{equation}
\end{theorem}
\begin{proof}
Theorem~\ref{thm:zeros} makes $\Norm_f$ nonzero on $\partial U$.
Near a zero of order $m$, write $\Norm_f(z)=(z-z_0)^m v(z)$ with
$v(z_0)\ne0$. Its logarithmic derivative is
$m/(z-z_0)+v'(z)/v(z)$, so the ordinary argument principle counts its zeros
with their complex multiplicities. These occur in conjugate pairs, except
for the real zeros, whose orders are even. Dividing the complex count by two
gives exactly \eqref{eq:multiplicity} and proves the formula.
\end{proof}
The integrand is scalar complex valued. A pointwise quaternionic expression
such as $f'(q)f(q)^{-1}$ cannot simply replace it in this proof.

\begin{theorem}[Rouch\'e on circular boundaries]
\label{thm:rouche}
Let $f,g$ be regular on $\Omega_D$. If
\begin{equation}
 |g(q)-f(q)|<|f(q)|\qquad(q\in\Gamma_U),
 \label{eq:rouche}
\end{equation}
then $\nu_g(U)=\nu_f(U)$.
\end{theorem}
\begin{proof}
For $0\le t\le1$, set $h_t=f+t(g-f)$. The strict inequality gives
$|h_t(q)|\ge|f(q)|-t|g(q)-f(q)|>0$ on $\Gamma_U$.
Consequently $\Norm_{h_t}$ has no zero on $\partial U$. The integral in
\eqref{eq:argument} depends continuously on $t$, since its integrand does and
its denominator remains uniformly separated from zero on the compact set
$[0,1]\times\partial U$. The integral is integer valued by
Theorem~\ref{thm:argument}, hence constant. Its endpoint values are the desired
counts.
\end{proof}
The hypothesis is imposed on \emph{all} quaternionic points of each boundary
sphere, not only on a single complex contour. This is the condition that
makes the zero-detection argument valid.

\begin{corollary}[A polynomial localization criterion]
If
\[
 P(q)=q^n+\sum_{k<n}q^k a_k,
 \qquad \sum_{k<n}R^k|a_k|<R^n,
\]
then all zeros of $P$ lie in $|q|<R$, with total multiplicity $n$.
In particular every zero satisfies
$|q|\le1+\max_{k<n}|a_k|$.
\end{corollary}
\begin{proof}
Apply Rouch\'e to $q^n$ and $P$ on the sphere $|q|=R$, and use
Corollary~\ref{cor:degreecount}. If $M=\max_{k<n}|a_k|$ and $R>1+M$, then
\[
 \sum_{k<n}R^k|a_k|\le M\frac{R^n-1}{R-1}<R^n.
\]
Let $R$ decrease to $1+M$. A nonmonic polynomial is first normalized by
right multiplication by its leading coefficient's inverse.
\end{proof}

\section{Reciprocals, minimum modulus, and open mappings}
\label{sec:reciprocal}
\subsection{The regular reciprocal}
Let $Z(f^s)$ be the union of all zero-bearing conjugacy classes of a nonzero
regular $f$. On $\Omega_D\setminus Z(f^s)$ define
\begin{equation}
 f^{-*}=(f^s)^{-1}f^c.
 \label{eq:reciprocal}
\end{equation}
The first factor is intrinsic and regular. Stem multiplication shows
$f*f^{-*}=f^{-*}*f=1$ on this domain. For a real-centered polynomial, for
example,
\[
 (q-p)^{-*}=\Delta_p(q)^{-1}(q-\bar p).
\]
The excluded set is the whole sphere $[p]$, even if $q-p$ has only the one
pointwise zero $p$.

\begin{proposition}[The conjugation change of variable]
\label{prop:T}
On $\Omega_D\setminus Z(f^s)$ put
\begin{equation}
 T_f(q)=f^c(q)^{-1}qf^c(q).
 \label{eq:T}
\end{equation}
This is a real-analytic diffeomorphism of that domain onto itself, preserving
each conjugacy class, and its inverse is $T_{f^c}$. Moreover,
\begin{align}
 f^{-*}(q)&=f(T_f(q))^{-1},\label{eq:inverse-eval}\\
 (f^{-*}*g)(q)&=f(T_f(q))^{-1}g(T_f(q)).
 \label{eq:quotient-eval}
\end{align}
\end{proposition}
\begin{proof}
Nonvanishing of $f^s(q)$ implies nonvanishing of $f^c(q)$ by the zero theorem.
Let $c=f^c(q)$, $r=T_f(q)$, and $n=f^s(q)$. Product evaluation applied to
$f^c*f$ gives
\[
 n=c f(r).
\]
Since $n$ is intrinsic, it commutes with $q$. Therefore
\[
 T_{f^c}(r)=f(r)^{-1}r f(r)
 =n^{-1}c(c^{-1}qc)c^{-1}n=n^{-1}qn=q.
\]
The same calculation in the opposite order proves the other inverse identity.
All maps are real analytic, and all preserve the same zero-free spheres.
Furthermore,
$n^{-1}c=(c f(r))^{-1}c=f(r)^{-1}$, proving
\eqref{eq:inverse-eval}.
Finally, product evaluation for $f^{-*}*g$ uses the conjugating value
$u=f^{-*}(q)=n^{-1}c$. Since $n$ commutes with $q$,
$u^{-1}qu=c^{-1}qc=T_f(q)$. This proves \eqref{eq:quotient-eval}.
\end{proof}
The reciprocal is therefore a pointwise inverse only \emph{after} a
sphere-preserving change of variable. This mechanism underlies the minimum
principle and geometric results developed in \cite{GStOpen,BS}.

\subsection{A minimum principle with no missing spheres}
\begin{theorem}[Minimum-modulus principle]
\label{thm:min}
Let $f\in\SR(\Omega_D)$. If $|f|$ has a local minimum at $p$, then either
$f(p)=0$ or $f$ is constant.
\end{theorem}
\begin{proof}
Assume $f(p)\ne0$. We first show that $[p]$ contains no zero of $f$.
This is immediate for real $p$. On a nonreal sphere write $f=A+IB$.
The squared norm
\[
 |A+IB|^2=|A|^2+|B|^2-2\Rea(A\bar B I)
\]
is an affine real function of the unit vector $I\in\SSS$.
A local minimum of an affine function on a sphere is a global minimum
(including the constant case). Thus the minimum at $p$ is no greater than
the value anywhere on $[p]$, which rules out a zero on that sphere.

The set $D'=D\setminus Z(\Norm_f)$ is connected: a polygonal path in a planar
domain can be perturbed around the finitely many points of a locally finite
set that meet a compact neighborhood of that path. It is conjugation
invariant and still contains a real interval, since real zeros are discrete.
Thus $f^{-*}$ is regular on the slice domain $\Omega_{D'}$.
By Proposition~\ref{prop:T}, $|f^{-*}|=1/|f\circ T_f|$ has a local maximum at
$T_f^{-1}(p)$. The maximum principle makes $f^{-*}$ constant, say equal to
$c\ne0$, on that domain. The identity $f*f^{-*}=1$ then gives
$f(q)c=1$ there. The identity principle extends this constant value of $f$
to all of $\Omega_D$.
\end{proof}
The initial sphere argument is essential: $f(p)\ne0$ alone does not guarantee
that the regular reciprocal is defined near $p$.

\subsection{The appropriate open-mapping statements}
Define the \emph{degenerate set}
\[
 D_f=\bigcup\{x+y\SSS:y>0,\ f\text{ is constant on }x+y\SSS\}.
\]
It is the union of spheres for which the coefficient $B(x,y)$ in
\eqref{eq:AB} vanishes.

\begin{theorem}[Open mapping with spherical qualifications]
\label{thm:open}
Let $f\in\SR(\Omega_D)$ be nonconstant.
\begin{enumerate}[label=\textup{(\roman*)}]
\item At every $p\notin D_f$, the image of every sufficiently small Euclidean
ball centered at $p$ contains a neighborhood of $f(p)$.
Consequently $f$ restricted to $\Omega_D\setminus\overline{D_f}$ is an open
map.
\item If $V\subseteq\Omega_D$ is open and axially symmetric, then $f(V)$ is
open. In particular $f(\Omega_D)$ is open.
\end{enumerate}
\end{theorem}
\begin{proof}
For (i), $h=f-f(p)$ has an isolated zero at $p$. Indeed, a real zero is
isolated; a nonreal zero is the unique zero on its sphere when that sphere is
not degenerate, and the zero-bearing classes are locally finite.
Choose a small closed ball $\overline{B(p,r)}$ in the domain whose boundary
contains no zero of $h$, and put
$m=\min_{\partial B(p,r)}|h|>0$.
If $|w-f(p)|<m/2$, then $|f-w|>m/2$ on the boundary but
$|f(p)-w|<m/2$. Hence $|f-w|$ attains its minimum in the interior. By
Theorem~\ref{thm:min}, that minimum is zero, since $f-w$ is nonconstant.
Thus $B(f(p),m/2)\subseteq f(B(p,r))$. Balls can also be chosen in
$\Omega_D\setminus\overline{D_f}$ when proving the restricted open-map claim.

For (ii), take $p\in V$. Choose a relatively compact axially symmetric
neighborhood $W$ of the entire class $[p]$ contained in $V$, with boundary
avoiding the zeros of $f-f(p)$. Such a neighborhood exists by the local
finiteness of zero-bearing classes: use a small disk around a real base point,
or two small conjugate disks around a nonreal pair, and circularize.
Now minimize $|f-w|$ on $\overline W$ and repeat the preceding argument.
It gives a neighborhood of $f(p)$ in $f(W)\subseteq f(V)$.
\end{proof}

\begin{example}[Why the unqualified open mapping theorem fails]
The entire regular function $f(q)=q^2$ is not open at $p=\ii$.
If $q^2=-1+t\jj$ with $t\ne0$ and $q=x+v$, then
$\Ima(q^2)=2xv=t\jj$. Thus $x\ne0$ and $v$ is parallel to $\jj$;
in particular $|q-\ii|\ge1$.
Consequently $f(B(\ii,1/2))$ contains $-1$ but none of the points
$-1+t\jj$, $t\ne0$, so it is not a neighborhood of $-1$.
The missing behavior occurs on a sphere that $f$ collapses to a point.
\end{example}
One must also retain the closure in the restricted open-map formulation.
For $q^2$, the degenerate spheres form the nonzero imaginary hyperplane;
their closure contains $0$. Removing only those nonreal spheres leaves $0$
but removes all preimages of negative real values, so the resulting
restricted image is not open at $0$.

\section{Schwarz--Pick and an intrinsic Riemann mapping theorem}
\label{sec:geometry}
\subsection{Regular fractional transformations}
For $a\in\BB$ define
\begin{equation}
 \mathcal M_a(q)=(1-q\bar a)^{-*}*(q-a).
 \label{eq:Mobius}
\end{equation}
The denominator has no pointwise zero in $\BB$, because $|q\bar a|<1$;
the zero theorem therefore makes its regular reciprocal well defined on the
whole ball. The two factors $1-q\bar a$ and $q-a$ commute under $*$ because
their coefficients belong to the same commutative plane containing $a$.

\begin{lemma}[Ball and boundary estimates]
\label{lem:Mobius}
The function $\mathcal M_a$ maps $\BB$ into $\BB$ and has exactly one zero,
at $a$. For $|a|<r<1$,
\begin{equation}
 \min_{|q|=r}|\mathcal M_a(q)|
 =\frac{r-|a|}{1-r|a|}.
 \label{eq:Mobius-min}
\end{equation}
\end{lemma}
\begin{proof}
Put $h(q)=1-q\bar a$. By \eqref{eq:quotient-eval},
\[
 \mathcal M_a(q)=(1-p\bar a)^{-1}(p-a),\qquad p=T_h(q).
\]
The elementary norm identity
\begin{equation}
 |1-p\bar a|^2-|p-a|^2=(1-|p|^2)(1-|a|^2)
 \label{eq:ballidentity}
\end{equation}
proves that its norm is less than one. Since $T_h$ is a bijection and
$h^c(a)=1-a^2$ commutes with $a$, $T_h(a)=a$, and the unique zero is $a$.
For $|p|=r$, write $t=\Rea(p\bar a)\in[-r|a|,r|a|]$. Then
\[
 \frac{|p-a|^2}{|1-p\bar a|^2}
 =\frac{r^2+|a|^2-2t}{1+r^2|a|^2-2t}.
\]
This ratio is decreasing in $t$, since its numerator minus denominator is
$-(1-r^2)(1-|a|^2)$. Its minimum is attained at $t=r|a|$, giving the square
of \eqref{eq:Mobius-min}. Such a $p$ exists, and $T_h$ permutes each sphere
and therefore the whole radius-$r$ three-sphere. The case $a=0$ gives the
same formula directly.
\end{proof}

\begin{theorem}[Schwarz lemma with an arbitrary zero]
\label{thm:schwarz-zero}
If $F:\BB\to\BB$ is regular and $F(a)=0$, then
\begin{equation}
 |F(q)|\le|\mathcal M_a(q)|\qquad(q\in\BB).
 \label{eq:schwarz-zero}
\end{equation}
Equality at any $q\ne a$ holds exactly when $F=\mathcal M_a u$ for a
constant $u$ of unit norm.
\end{theorem}
\begin{proof}
By regular division, $F=(q-a)*G$. Define
\[
 H=(1-q\bar a)*G.
\]
It is regular throughout $\BB$ and $F=\mathcal M_a*H$.
Off the zero-bearing sphere of $\mathcal M_a$,
$H=\mathcal M_a^{-*}*F$. For $r>|a|$, quotient evaluation and
\eqref{eq:Mobius-min} therefore give, on $|q|=r$,
\[
 |H(q)|\le\frac{1-r|a|}{r-|a|}.
\]
Apply maximum modulus on $B(0,r)$ and let $r\uparrow1$ to obtain $|H|\le1$
on $\BB$. Evaluating $F=\mathcal M_a*H$ with
Proposition~\ref{prop:evaluation} yields \eqref{eq:schwarz-zero}.
If equality holds where $\mathcal M_a(q)\ne0$, then $|H|=1$ at the
corresponding conjugate point. Maximum modulus makes $H$ a constant $u$
with $|u|=1$. The converse follows by direct evaluation.
\end{proof}

\begin{theorem}[Regular Schwarz--Pick inequality]
\label{thm:pick}
Let $f:\BB\to\BB$ be regular, $a\in\BB$, and $b=f(a)$. Define
\begin{equation}
 \Phi_{f,b}=(1-f\bar b)^{-*}*(f-b).
 \label{eq:normalized-f}
\end{equation}
Then
\begin{equation}
 |\Phi_{f,b}(q)|\le|\mathcal M_a(q)|\qquad(q\in\BB).
 \label{eq:pick}
\end{equation}
Equality at some $q\ne a$ holds exactly when
$\Phi_{f,b}=\mathcal M_a u$ for a constant $|u|=1$.
In particular, if $a$ is real,
\begin{equation}
 \frac{|\pc f(a)|}{1-|f(a)|^2}\le\frac1{1-a^2}.
 \label{eq:real-pick-derivative}
\end{equation}
\end{theorem}
\begin{proof}
The function $h=1-f\bar b$ is regular and zero free on $\BB$, since
$|f(q)||b|<1$. Its reciprocal is therefore regular everywhere in the ball.
Quotient evaluation and \eqref{eq:ballidentity}, now with
$p=f(T_h(q))$ and $a=b$, show $|\Phi_{f,b}(q)|<1$.
The identity $h*\Phi_{f,b}=f-b$ evaluated at $a$ gives
\[
 (1-|b|^2)\Phi_{f,b}(a)=0,
\]
because $h(a)=1-|b|^2$ is a real nonzero scalar and hence does not rotate
the argument. Theorem~\ref{thm:schwarz-zero} applies to $\Phi_{f,b}$ and
proves the norm inequality and its equality condition.

If $a$ is real, restrict the identity $h*\Phi_{f,b}=f-b$ to the real axis
and differentiate at $a$. Since $\Phi_{f,b}(a)=0$,
\[
 \pc\Phi_{f,b}(a)=\frac{\pc f(a)}{1-|b|^2}.
\]
Also $\mathcal M_a(x)=(x-a)/(1-ax)$ on the real axis and
$\mathcal M_a'(a)=1/(1-a^2)$. Divide \eqref{eq:pick} at $q=x$ by $|x-a|$
and let real $x\to a$.
\end{proof}
This is the regular-quotient version of the quaternionic Schwarz--Pick
phenomenon in \cite{BS}. It is not obtained by merely replacing the complex
variables in the ordinary pointwise fraction. At nonreal base points, the
real differential has the two pieces in \eqref{eq:differential}; the simple
real-base differential inequality above is deliberately not asserted there.

For comparison with a common alternative notation, the algebraic identity
\[
 (1-f\bar b)*(f-b)=(f-b)*(1-\bar b*f)
\]
shows that
\[
 (1-f\bar b)^{-*}*(f-b)=(f-b)*(1-\bar b*f)^{-*}.
\]
Here $\bar b*f$ means the regular product, not ordinary left multiplication of
function values. For an equality case, one can recover $f$ from
$F=\mathcal M_a u$ by the explicit regular fraction
\[
 f=(1+F\bar b)^{-*}*(F+b).
\]
Indeed, expanding either side after multiplication by its denominator reduces
to the preceding algebraic identity with a plus sign. The denominator is zero
free because $|F||b|<1$, and \eqref{eq:ballidentity} with $-b$ shows that this
recovery map again takes values in the ball.

\subsection{A genuine Riemann mapping theorem in the intrinsic category}
Ordinary composition of regular functions need not be regular. For example,
$q\mapsto q\ii$ and $q\mapsto q^2$ are regular, but their composite
$q\mapsto q\ii q\ii$ equals $-|q|^2$ on $\C_{\jj}$, which is not
holomorphic. Composition does work when the inner function is intrinsic:
if $\varphi$ has a scalar holomorphic stem $\phi$, then the stem of
$f\circ\varphi$ is $\mathcal F\circ\phi$, provided the image lies in the
appropriate domain.

\begin{theorem}[Intrinsic Riemann mapping theorem]
\label{thm:riemann}
Let $D\subsetneq\C$ be simply connected, open, conjugation invariant, and
let $c\in D\cap\R$. There is a unique intrinsic slice-regular bijection
\[
 \varphi:\Omega_D\longrightarrow\BB
\]
with $\varphi(c)=0$ and $\pc\varphi(c)>0$ real. Its inverse is also intrinsic
and slice regular.
\end{theorem}
\begin{proof}
The classical normalized Riemann mapping theorem provides a unique
biholomorphism $\phi:D\to\{z:|z|<1\}$ with $\phi(c)=0$ and
$\phi'(c)>0$. The function $z\mapsto\overline{\phi(\bar z)}$ has the same
normalization. Uniqueness therefore gives
$\phi(\bar z)=\overline{\phi(z)}$.
Its inverse $\psi$ has the same symmetry. Both functions consequently have
intrinsic slice extensions $\varphi$ and $\chi$.

On every $\C_I$, the extensions are precisely the scalar maps $\phi$ and
$\psi$ under the identification $x+\iota y\leftrightarrow x+Iy$. Hence
$\chi\circ\varphi=\Id$ and $\varphi\circ\chi=\Id$ section by section,
proving both bijectivity and regularity of the inverse.
For uniqueness, any intrinsic bijection satisfying the normalization restricts
to a biholomorphism on each section: a nonreal point cannot map to a real
value, since its entire sphere would then have the same image and injectivity
would fail. Surjectivity and preservation of sections thus give bijectivity of
each section restriction. Apply normalized scalar uniqueness on one section
and the representation formula.
\end{proof}
The theorem asserts a slice-analytic equivalence, not arbitrary conformal
uniformization in $\R^4$. For an intrinsic function induced by
$\phi=u+\iota v$, the two local stretch factors in
\eqref{eq:differential} are $|\phi'(z)|$ and $|v(x,y)/y|$. They can differ:
for $q^2$ at $1+\ii$ they are $2\sqrt2$ and $2$.

\section{Laurent expansions, residues, and essential singularities}
\label{sec:laurent}
The statements in this section concern isolated singularities at \emph{real}
points. This restriction makes ordinary radial annuli and ordinary integer
powers the right objects. Nonreal singularities require keeping track of an
entire conjugacy sphere and lead to a more elaborate theory; see
\cite{GPStSing}.

\subsection{Laurent's theorem}
\begin{theorem}[Laurent expansion on a spherical annulus]
\label{thm:laurent}
Let $f$ be slice regular on $A_{r,R}=\{q:r<|q|<R\}$, where
$0\le r<R\le\infty$. Then
\begin{equation}
 f(q)=\sum_{n\in\mathbb Z}q^n a_n
 \label{eq:laurent}
\end{equation}
with unique quaternionic coefficients, converging absolutely and uniformly
on compact subannuli. For $r<\rho<R$,
\begin{equation}
 a_n=\frac1{2\pi}\int_{|s|=\rho,\ s\in\C_I}
        s^{-n-1}\,\mathrm ds_I\,f(s),
 \qquad |a_n|\le M_I(\rho)\rho^{-n}.
 \label{eq:laurentcoef}
\end{equation}
The coefficient is independent of $\rho$ and $I$.
\end{theorem}
\begin{proof}
The stem $\mathcal F$ is holomorphic on the complex annulus. Apply the scalar
Laurent theorem to its four complex coordinates. Its symmetry
$\mathcal F(\bar z)=\kappa(\mathcal F(z))$, together with uniqueness of the
Laurent coefficients, implies that each coefficient is fixed by $\kappa$,
hence belongs to $\HH$. Evaluating the resulting stem series gives
\eqref{eq:laurent}. The complex integral formulas, or the splitting lemma,
give \eqref{eq:laurentcoef}. Coefficient estimates on a radius smaller than
the compact subannulus control the negative tail, and estimates on a larger
radius control the positive tail, both by geometric series. This proves the
stated absolute and uniform convergence.
\end{proof}
Translation by a real center $c$ replaces $q^n$ by $(q-c)^n$ without changing
the argument.

\begin{corollary}[Classification and Riemann removability at a real point]
\label{cor:removable-slice}
Let $f$ be regular on $0<|q|<R$. Its singularity at $0$ is removable if
all negative Laurent coefficients vanish, a pole if finitely many but at
least one are nonzero, and essential if infinitely many are nonzero.
If $f$ is bounded near zero, the singularity is removable. If the pole has
order $m$, then
\begin{equation}
 |f(q)|=|q|^{-m}\bigl(|a_{-m}|+o(1)\bigr)\quad(q\to0).
 \label{eq:polegrowth}
\end{equation}
\end{corollary}
\begin{proof}
The Laurent theorem gives the trichotomy. Under a local bound $M$,
\eqref{eq:laurentcoef} gives $|a_{-m}|\le M\rho^m$ for every small $\rho$,
so each negative coefficient vanishes as $\rho\to0$.
For a pole, $q^m f(q)$ extends regularly with nonzero value $a_{-m}$ at
zero; continuity and multiplicativity of the norm prove \eqref{eq:polegrowth}.
\end{proof}

\begin{theorem}[Residue theorem for isolated real singularities]
\label{thm:residue-slice}
Suppose $f$ is regular on a neighborhood of
$\overline{\Omega_U}\setminus\{c_1,\ldots,c_N\}$, where the $c_j$ are
real interior points and $U$ has piecewise smooth positively oriented
boundary. Define $\Res_{c_j}f$ as the coefficient of $(q-c_j)^{-1}$ in its
Laurent expansion. Then
\begin{equation}
 \frac1{2\pi}\int_{\partial U_I}\mathrm ds_I\,f(s)
     =\sum_{j=1}^N\Res_{c_j}f.
 \label{eq:residue-slice}
\end{equation}
\end{theorem}
\begin{proof}
Remove disjoint small disks around the $c_j$ in the $I$-section and apply
Cauchy's theorem to both splitting components. On each small circle,
uniform convergence of the Laurent expansion permits integration term by
term. All powers integrate to zero except $(s-c_j)^{-1}$, whose normalized
integral is $1$. The right coefficients remain on the right. The Laurent
coefficient's independence of $I$ gives the stated quaternionic formula.
\end{proof}

\subsection{Casorati--Weierstrass}
\begin{theorem}[Density near a real essential singularity]
\label{thm:casorati}
If $f$ has an essential singularity at the real point $0$, then for every
$0<r<R$ the set $f(\{0<|q|<r\})$ is dense in $\HH$.
\end{theorem}
\begin{proof}
Suppose an open ball $B(a,\varepsilon)$ is omitted by the image of a punctured
ball. Then $g=f-a$ has no zero there and $|g|\ge\varepsilon$.
Its regular reciprocal satisfies
\[
 |g^{-*}(q)|=|g(T_g(q))|^{-1}\le\varepsilon^{-1},
\]
because $T_g$ preserves each radius. By removability, $h=g^{-*}$ extends
regularly across zero. It is not identically zero, since its regular product
with $g$ is $1$ off zero.
If $h(0)\ne0$, then $h^{-*}$ is regular near zero and agrees with $g$
there off zero, so $f$ has a removable singularity.
If $h(0)=0$, write $h=q^m u$ with $m\ge1$ and $u(0)\ne0$ using its Taylor
series. Near zero, $u^{-*}$ is regular and
\[
 g=h^{-*}=q^{-m}u^{-*}.
\]
Its leading coefficient is nonzero, so $g$, and hence $f$, has a pole.
Both conclusions contradict essentiality.
\end{proof}
The proof is the familiar reciprocal argument from complex analysis, but
only after the regular reciprocal and its conjugation change of variable
have been established.

\section{Normal families and Hurwitz}
\label{sec:normal}
\begin{theorem}[Weierstrass convergence and Montel]
\label{thm:montel}
A locally uniform limit of slice-regular functions on $\Omega_D$ is slice
regular. Every family of slice-regular functions that is bounded uniformly
on each compact subset of $\Omega_D$ is normal: every sequence has a
subsequence converging locally uniformly to a slice-regular function.
\end{theorem}
\begin{proof}
Fix orthogonal $I,J$ and split each restriction to $\C_I$ into two complex
components. For the first assertion, scalar Weierstrass convergence gives
holomorphic limit components. Their slice extension is regular, and the
representation formula shows that it equals the given limit.

For the second assertion, scalar Montel applied to both components, with a
successive subsequence extraction, gives a subsequence converging uniformly
on compact subsets of the chosen section. To propagate convergence to a
compact $K\subset\Omega_D$, use its compact two-point section projection
\[
 K_I=\{\Rea q\pm I|\Ima q|:q\in K\}.
\]
The coefficients in \eqref{eq:representation} each have norm at most one.
Thus the uniform difference of two members on $K$ is at most twice their
uniform difference on $K_I$. The subsequence therefore converges on every
compact subset of $\Omega_D$, and the first assertion proves regularity of
the limit.
\end{proof}

\begin{theorem}[Hurwitz for quaternionic zero-free functions]
\label{thm:hurwitz}
Suppose $f_n\in\SR(\Omega_D)$ have no zeros in $\Omega_D$ and converge
locally uniformly to $f$. Then either $f\equiv0$ or $f$ is zero free.
\end{theorem}
\begin{proof}
The representation formula recovers the stem components $A_n,B_n$ by taking
linear combinations of values on a fixed section. Hence their convergence
is locally uniform, and \eqref{eq:normstem} gives
$\Norm_{f_n}\to\Norm_f$ locally uniformly on $D$.
Theorem~\ref{thm:zeros} says that every $\Norm_{f_n}$ is zero free.
Scalar Hurwitz implies that $\Norm_f$ is zero free or identically zero.
The first case makes $f$ zero free; in the second, \eqref{eq:normonreal} and
the identity principle imply $f\equiv0$.
\end{proof}
Merely viewing $f_I$ as a holomorphic map into $\C^2$ would not justify this
last theorem: the maps $z\mapsto(z,1/n)$ are nonvanishing but converge to
$(z,0)$, which has a zero. Quaternionic zero-freeness on \emph{all} spheres,
encoded by the scalar norm stem, is the additional information that makes
Hurwitz possible.

\section{The Cauchy--Fueter operator and the four-dimensional Cauchy formula}
\label{sec:fueter-cauchy}
We now change the notion of regularity. Let $\Omega\subset\HH\cong\R^4$ be
an arbitrary connected open set; no axial symmetry is assumed. Put
\begin{align}
 \cD&=\partial_{x_0}+\ii\partial_{x_1}
             +\jj\partial_{x_2}+\kk\partial_{x_3},\\
 \overline{\cD}&=\partial_{x_0}-\ii\partial_{x_1}
             -\jj\partial_{x_2}-\kk\partial_{x_3}.
\end{align}
A $C^1$ function $f:\Omega\to\HH$ is \emph{left Cauchy--Fueter regular},
or left monogenic, when $\cD f=0$. For a function on the other side of a
product, use the right-acting operator
\[
 g\cD=\partial_{x_0}g+(\partial_{x_1}g)\ii
                  +(\partial_{x_2}g)\jj+(\partial_{x_3}g)\kk.
\]
Anticommutation of the imaginary units gives
\begin{equation}\label{eq:factor-laplacian}
 \overline{\cD}\cD=\cD\overline{\cD}=\lap,
 \qquad \lap=\sum_{\mu=0}^{3}\partial_{x_\mu}^2.
\end{equation}
These are identities of constant-coefficient differential operators. Notice
at once that $\cD q=1-1-1-1=-2$: the identity function, a basic example in
slice analysis, is not Fueter regular. The integral theory developed next is
classical; Sudbery~\cite{Sudbery} is a particularly useful reference.

\begin{lemma}[Quaternionic Green identity]\label{lem:green}
Let $\Omega$ be bounded with $C^1$ boundary, let $n=n_0+\ii n_1+\jj n_2+\kk n_3$
be its outward unit normal, and let $f,g$ be $C^1$ near $\overline\Omega$.
Then
\begin{equation}\label{eq:green}
 \int_{\partial\Omega}g(y)n(y)f(y)\dd S_y
 =\int_\Omega\bigl((g\cD)(y)f(y)+g(y)(\cD f)(y)\bigr)\dd V_y.
\end{equation}
\end{lemma}
\begin{proof}
Write $e_0=1,e_1=\ii,e_2=\jj,e_3=\kk$. Apply the ordinary divergence theorem,
component by component, to the four-vector with quaternionic components
$g e_\mu f$. Its divergence is
\[
 \sum_\mu\partial_{x_\mu}(g e_\mu f)
 =\sum_\mu(\partial_{x_\mu}g)e_\mu f
       +\sum_\mu g e_\mu\partial_{x_\mu}f
 =(g\cD)f+g(\cD f).
\]
The order of the three factors in the boundary integral is thereby fixed.
\end{proof}
In particular, a Fueter-regular function satisfies the Cauchy theorem
\begin{equation}\label{eq:fueter-cauchy-theorem}
 \int_{\partial\Omega}n(y)f(y)\dd S_y=0.
\end{equation}
The integration cycle is now three-dimensional, not a planar curve.

Define the fundamental kernel by
\begin{equation}\label{eq:fueter-kernel}
 E(q)=\frac{\bar q}{2\pi^2|q|^4},\qquad q\ne0.
\end{equation}
It is both left and right Fueter regular away from zero. For example,
$\cD\bar q=4$ and
\[
 \sum_{\mu=0}^3 e_\mu\bar q\,\partial_{x_\mu}|q|^{-4}
 =-4q\bar q\,|q|^{-6}=-4|q|^{-4},
\]
which cancels the derivative of $\bar q$; the right calculation is identical
with the order reversed. On $|q|=r$, the outward normal is $q/r$, so
\begin{equation}\label{eq:kernel-flux}
 E(q)n(q)=n(q)E(q)=\frac{1}{2\pi^2r^3},
 \qquad \int_{|q|=r}E(q)n(q)\dd S_q=1,
\end{equation}
since the unit three-sphere has area $2\pi^2$.

\begin{theorem}[Cauchy--Pompeiu and Cauchy--Fueter formulas]
\label{thm:fueter-cauchy}
Under the boundary hypotheses of Lemma~\ref{lem:green}, for $f\in C^1$
near $\overline\Omega$ and $x\in\Omega$,
\begin{equation}\label{eq:pompeiu}
 f(x)=\int_{\partial\Omega}E(y-x)n(y)f(y)\dd S_y
       -\int_\Omega E(y-x)(\cD f)(y)\dd V_y.
\end{equation}
Consequently, if $\cD f=0$, then
\begin{equation}\label{eq:fueter-cauchy}
 \boxed{\displaystyle
 f(x)=\int_{\partial\Omega}E(y-x)n(y)f(y)\dd S_y.}
\end{equation}
\end{theorem}
\begin{proof}
Apply \eqref{eq:green} with $g(y)=E(y-x)$ to
$\Omega\setminus\overline{B(x,\varepsilon)}$. The right derivative
$g\cD$ vanishes there. The normal of the punctured domain on its inner
boundary is the negative of the outward normal of the small ball. Moving
that boundary term to the other side gives
\[
 \int_{|y-x|=\varepsilon}E(y-x)n_{\rm ball}(y)f(y)\dd S_y
 =\int_{\partial\Omega}E(y-x)n(y)f(y)\dd S_y
   -\int_{\Omega\setminus\overline{B(x,\varepsilon)}}
          E(y-x)(\cD f)(y)\dd V_y.
\]
The left side tends to $f(x)$ by \eqref{eq:kernel-flux} and continuity. The
volume integral converges because $|E(y-x)|=O(|y-x|^{-3})$, which is locally
integrable in four dimensions. This proves both formulas.
\end{proof}
Equation~\eqref{eq:fueter-cauchy} represents functions already known to be
regular; it is not a claim that arbitrary boundary values admit a regular
extension.

\section{Fueter analogues of the basic analytic theorems}
\label{sec:fueter-consequences}
\subsection{Analyticity, Cauchy estimates, and Liouville's theorem}
\begin{proposition}\label{prop:fueter-analytic}
Every $C^1$ Fueter-regular function is real analytic, and each of its four
real components is harmonic.
\end{proposition}
\begin{proof}
Fix a closed ball strictly inside the domain and use
\eqref{eq:fueter-cauchy} on that ball. For interior $x$, the kernel is smooth
in $x$ uniformly for $y$ on the boundary, so all derivatives pass under the
integral. To obtain real analyticity near the center $a$, expand
\[
 |y-x|^{-4}
 =|y-a|^{-4}\left(1+
 \frac{-2\langle y-a,x-a\rangle+|x-a|^2}{|y-a|^2}\right)^{-2}.
\]
For $|x-a|$ sufficiently small compared with the ball radius, the binomial
series converges absolutely and uniformly on the integration sphere. After
multiplication by $\overline{y-x}$ it is a convergent power series in the
four real coordinates of $x-a$, and it can be integrated term by term.
Finally \eqref{eq:factor-laplacian} gives
$\lap f=\overline{\cD}(\cD f)=0$.
\end{proof}

\begin{theorem}[Cauchy derivative estimates]\label{thm:fueter-estimates}
Suppose $f$ is Fueter regular near $\overline{B(a,R)}$, and let
$M_R=\max_{|q-a|=R}|f(q)|$. For every multi-index $\alpha$,
\begin{equation}\label{eq:fueter-estimate}
 |\partial^\alpha f(a)|\le
 C_\alpha\frac{M_R}{R^{|\alpha|}},
 \qquad
 C_\alpha=2\pi^2\max_{|\omega|=1}|\partial^\alpha E(\omega)|.
\end{equation}
In particular $C_0=1$, and each first partial derivative satisfies the
estimate with $C_\alpha=5$.
\end{theorem}
\begin{proof}
Differentiate \eqref{eq:fueter-cauchy}. Since $E$ is homogeneous of degree
$-3$, its derivative $\partial^\alpha E$ is homogeneous of degree
$-3-|\alpha|$. Taking norms and using the area $2\pi^2R^3$ proves
\eqref{eq:fueter-estimate}. Direct differentiation gives
\[
 \partial_{x_\mu}E(q)=\frac{1}{2\pi^2}
 \left(\bar e_\mu|q|^{-4}-4x_\mu\bar q|q|^{-6}\right),
\]
whose norm is at most $5/(2\pi^2|q|^4)$. This constant is sufficient; no
sharpness is asserted for it.
\end{proof}

\begin{corollary}[Liouville and polynomial growth]\label{cor:fueter-liouville}
A bounded entire Fueter-regular function is constant. More generally, if
\[
 |f(q)|\le C(1+|q|)^\beta\qquad(q\in\HH),\quad\beta\ge0,
\]
then $f$ is a polynomial in the four real coordinates, of total degree at
most $\lfloor\beta\rfloor$, satisfying $\cD f=0$.
\end{corollary}
\begin{proof}
For bounded $f$, let $R\to\infty$ in the first-derivative estimates at an
arbitrary center. All first partials vanish. Under the growth hypothesis,
$M_R\le C(1+|a|+R)^\beta$, so every derivative of order strictly greater than
$\beta$ vanishes at every center. The real Taylor expansion therefore
terminates with the asserted degree.
\end{proof}
The polynomial conclusion differs from the slice conclusion: here a
polynomial is a polynomial in $x_0,x_1,x_2,x_3$, not necessarily a polynomial
in the single quaternionic variable $q$ with right coefficients.

\subsection{Mean values and maximum modulus}
\begin{theorem}[Mean value and maximum modulus]\label{thm:fueter-max}
If $f$ is Fueter regular near $\overline{B(a,R)}$, then
\begin{equation}\label{eq:fueter-mean}
 f(a)=\frac{1}{2\pi^2R^3}\int_{|y-a|=R}f(y)\dd S_y.
\end{equation}
If $|f|$ has a local maximum in a connected domain, then $f$ is constant.
\end{theorem}
\begin{proof}
Formula \eqref{eq:fueter-mean} follows immediately from
\eqref{eq:fueter-cauchy} and \eqref{eq:kernel-flux}. Suppose $|f|\le M=|f(a)|$
near $a$. The case $M=0$ gives local constancy immediately. Otherwise the
real-linear functional
\[
 \ell(w)=\frac{\Rea(w\overline{f(a)})}{M}
\]
satisfies $\ell(w)\le |w|$. On every sufficiently small sphere centered at
$a$, the average of $\ell(f)$ is $M$, while $\ell(f)\le |f|\le M$.
Continuity forces equality at every point of the sphere. Equality in the
real Euclidean Cauchy--Schwarz inequality, together with $|f|\le M$, then
forces $f=f(a)$ there. Varying the radius proves constancy on a ball.
Real analyticity propagates that constancy through the connected domain.
\end{proof}

\subsection{A flux version of Morera's theorem}
\begin{theorem}[Fueter--Morera]\label{thm:fueter-morera}
A continuous function $f:\Omega\to\HH$ is Fueter regular if and only if
\begin{equation}\label{eq:fueter-morera}
 \int_{\partial B}n(y)f(y)\dd S_y=0
\end{equation}
for every ball whose closure lies in $\Omega$.
\end{theorem}
\begin{proof}
Necessity is \eqref{eq:fueter-cauchy-theorem}. For sufficiency, work in a
relatively compact subdomain and convolve $f$ with a real smooth mollifier,
obtaining $f_\varepsilon$. By Fubini's theorem, the flux of
$f_\varepsilon$ through a ball is the average of the fluxes of $f$ through
translated balls, hence is zero when the translations stay in $\Omega$.
The divergence theorem gives
\[
 \int_B\cD f_\varepsilon\dd V=0
\]
for every such ball. Shrinking balls to their centers proves
$\cD f_\varepsilon=0$ pointwise. On a fixed smaller ball, apply the Cauchy
formula to $f_\varepsilon$ and pass to the limit, using uniform convergence
on the boundary. The limiting integral is smooth and left regular inside
because $\cD_x E(y-x)=0$ there. It equals $f$, so $f$ is regular locally and
therefore throughout $\Omega$.
\end{proof}

\subsection{Removable singularities and residues}
\begin{theorem}[A sharp growth threshold for removability]
\label{thm:fueter-removable}
Suppose $f$ is Fueter regular on $B(a,R)\setminus\{a\}$ and
\begin{equation}\label{eq:fueter-removable-growth}
 \sup_{|q-a|=r}|f(q)|=o(r^{-3})\qquad(r\downarrow0).
\end{equation}
Then $f$ extends uniquely to a Fueter-regular function on $B(a,R)$.
In particular a bounded isolated singularity is removable. The exponent
$3$ cannot be replaced by a statement allowing all $O(r^{-3})$ growth.
\end{theorem}
\begin{proof}
Choose $R_0<R$. Apply the Cauchy formula on the punctured ball with outer
radius $R_0$ and inner radius $\varepsilon$, at a fixed point $x\ne a$.
For sufficiently small $\varepsilon$, the kernel $E(y-x)$ is bounded on the
inner sphere. The norm of the inner boundary contribution is therefore at
most a constant depending on $x$ times
\[
 \varepsilon^3\sup_{|y-a|=\varepsilon}|f(y)|,
\]
which tends to zero. Thus $f(x)$ equals the outer boundary integral alone.
That integral is regular throughout $B(a,R_0)$, including at $a$. The
extensions agree on overlaps by continuity. Finally $E(q-a)c$, with
$c\ne0$, is a nonremovable example of exact order $|q-a|^{-3}$.
\end{proof}

\begin{definition}
For a Fueter-regular function with an isolated singularity at $a$, define
its \emph{Fueter residue} by
\begin{equation}\label{eq:fueter-residue}
 \Res^{F}_{a}f=\int_{|q-a|=r}n(q)f(q)\dd S_q,
\end{equation}
where $r$ is any sufficiently small positive radius and the normal points
away from $a$. This normalization includes the $2\pi^2$ in the kernel,
not in the definition of the residue.
\end{definition}
The value is independent of $r$, by applying the Green identity to an
annulus. This elementary flux residue need not encode the whole principal
part of a singularity.

\begin{theorem}[Fueter residue theorem]\label{thm:fueter-residue}
Let $\Omega$ be a bounded $C^1$ domain. Suppose $f$ is Fueter regular near
$\overline\Omega$ except at finitely many
interior points $a_1,\ldots,a_m$, where it is regular in punctured
neighborhoods. Then
\begin{equation}\label{eq:fueter-residue-theorem}
 \int_{\partial\Omega}n(q)f(q)\dd S_q
       =\sum_{j=1}^m\Res^F_{a_j}f.
\end{equation}
In particular $\Res^F_a(E(q-a)c)=c$.
\end{theorem}
\begin{proof}
Remove pairwise disjoint small balls about the singularities and apply
\eqref{eq:fueter-cauchy-theorem} to the remaining domain. Its inner normals
are negatives of the normals in \eqref{eq:fueter-residue}; moving those
terms to the other side proves the formula. The final assertion follows
from $nE=1/(2\pi^2r^3)$ on the sphere.
\end{proof}
A zero residue does not imply removability. The function $\partial_{x_0}E$
is regular away from zero, is nonzero and homogeneous of degree $-4$, and
is singular at zero. Since $E$ is odd, $\partial_{x_0}E$ is even, whereas
$n$ is odd on a centered sphere. Its flux, and hence its residue, is zero.

\subsection{Two complex theorems that fail without modification}
\begin{example}\label{ex:fueter-plane}
The function
\begin{equation}\label{eq:fueter-plane}
 L(q)=x_1-\ii x_0
\end{equation}
is Fueter regular because $\cD L=-\ii+\ii=0$. It vanishes on the entire
plane $\{x_0=x_1=0\}$ but is not identically zero. Its image is contained in
$\C_\ii$ and thus has empty interior in $\HH$. Therefore neither the
complex accumulation-point identity theorem nor the nonconstant open
mapping theorem holds for arbitrary Fueter-regular functions. The
appropriate elementary uniqueness statement here is vanishing on a
nonempty open set, which does imply global vanishing by real analyticity.
\end{example}

\section{Fueter's mapping theorem and a constructive local inverse}
\label{sec:fueter-map}
The two theories are not unrelated alternatives: a second-order
differential operator connects them. We give the quaternionic version of
Fueter's mapping theorem directly, followed by an elementary axial inverse.
The broader inverse theory is developed in~\cite{CSSInverse,CPSSInverse}.

\subsection{Axial differential identities}
Away from the real axis write $q=x+Ir$, where $r=|\Ima q|>0$, and consider
an axial function
\[
 f(q)=A(x,r)+I B(x,r),
\]
with quaternion-valued $A,B$. Direct differentiation gives
\begin{align}
 \cD f
   &=\left(A_x-B_r-\frac{2B}{r}\right)+I(B_x+A_r),
      \label{eq:axial-D}\\
 \lap f
   &=A_{xx}+A_{rr}+\frac{2A_r}{r}
      +I\left(B_{xx}+B_{rr}+\frac{2B_r}{r}-\frac{2B}{r^2}\right).
      \label{eq:axial-lap}
\end{align}
For completeness, put $v=\Ima q=\ii x_1+\jj x_2+\kk x_3$, so $I=v/r$.
For $\ell=1,2,3$,
\[
 \partial_{x_\ell}r=\frac{x_\ell}{r},\qquad
 \partial_{x_\ell}I=\frac{e_\ell}{r}-\frac{v x_\ell}{r^3}.
\]
Consequently $\sum_{\ell=1}^3e_\ell\partial_{x_\ell}I=-2/r$,
$\Delta_{\R^3}I=-2I/r^2$, and
$\sum_\ell(\partial_{x_\ell}I)(\partial_{x_\ell}r)=0$.
The product rule and the ordinary radial Laplacian in three dimensions now
give \eqref{eq:axial-D}--\eqref{eq:axial-lap}, with all coefficients on the
indicated sides.

For a slice-regular function the stem Cauchy--Riemann equations imply
$A_x=B_r$, $A_r=-B_x$, and $A_{xx}+A_{rr}=B_{xx}+B_{rr}=0$. Hence
\begin{equation}\label{eq:slice-D-lap}
 \cD f=-\frac{2B}{r},\qquad
 \lap f=\frac{2A_r}{r}
       +I\left(\frac{2B_r}{r}-\frac{2B}{r^2}\right).
\end{equation}

\begin{theorem}[Fueter mapping, biharmonicity, and kernel]
\label{thm:fueter-map}
If $f\in\SR(\Omega_D)$, then
\begin{equation}\label{eq:fueter-map}
 \boxed{\cD(\lap f)=0,\qquad \lap^2 f=0.}
\end{equation}
On this connected slice domain the kernel of $f\mapsto\lap f$ consists
exactly of the affine functions $q\mapsto qa+b$, with $a,b\in\HH$.
Moreover, a function that is both slice regular and left Fueter regular
is constant.
\end{theorem}
\begin{proof}
Away from the real axis put
\[
 C=\frac{2A_r}{r},\qquad H=\frac{2B_r}{r}-\frac{2B}{r^2}.
\]
The Cauchy--Riemann equations give
\begin{align*}
 C_x&=\frac{2B_{rr}}r=H_r+\frac{2H}r,\\
 H_x+C_r
 &=\frac{2(B_{rx}+A_{rr})}{r}
     -\frac{2(B_x+A_r)}{r^2}=0.
\end{align*}
Equation~\eqref{eq:axial-D} therefore says that
$\cD(C+IH)=\cD(\lap f)=0$. Applying $\overline{\cD}$ proves
$\lap^2 f=0$. Slice-regular functions are smooth across the real axis, as
established by the real-centered Taylor theorem, so continuity extends
both differential identities across it.

If $\lap f=0$, comparison of the values at $I$ and $-I$ gives $C=H=0$.
Thus $A_r=0$ and $rB_r=B$. Locally in the upper half-plane,
$B=r b(x)$, and the Cauchy--Riemann equations then give $b'(x)=0$ and
$A_x=b$. Hence $f(q)=qa+b_0$ on a nonempty open set, for fixed
quaternions $a,b_0$. The identity theorem extends this formula through
$\Omega_D$. Conversely affine functions have zero Laplacian.
Finally $\cD f=0$ in \eqref{eq:slice-D-lap} forces $B=0$; the
Cauchy--Riemann equations force $A$ to be locally constant, and the identity
theorem makes it globally constant.
\end{proof}

\begin{example}[The two Cauchy kernels are linked]\label{ex:kernel-link}
Writing $v=\Ima q$ and $x=x_0$, direct calculation gives
\begin{align}
 \lap 1&=\lap q=0,&
 \lap q^2&=-4,\\
 \lap q^3&=-4(3x+v),&
 \lap(q^{-1})&=-\frac{4\bar q}{|q|^4}=-8\pi^2E(q),\quad q\ne0.
 \label{eq:kernel-link}
\end{align}
For the last identity, write $q^{-1}=\bar q\,|q|^{-2}$. In four dimensions
$\lap|q|^{-2}=0$ away from zero, and the cross term in the product rule is
$2\sum_\mu\bar e_\mu\partial_{x_\mu}|q|^{-2}
=-4\bar q|q|^{-4}$. Thus the slice kernel of homogeneity $-1$ becomes,
under the Laplacian, the Fueter kernel of homogeneity $-3$.
\end{example}

\subsection{An inverse obtained by two integrations}
A local inverse should only be expected within the axial class. In the
following theorem the planar set lies strictly above the real axis. Its
circularization is therefore a \emph{product domain}, rather than a slice
domain meeting the real line. We use the stem definition there; the global
identity and maximum principles proved earlier are not being asserted for
arbitrary product domains.

\begin{theorem}[Constructive axial inverse]\label{thm:fueter-inverse}
Let $U\subset\{(x,r)\in\R^2:r>0\}$ be nonempty, simply connected and open.
Suppose
\[
 g(x+Ir)=P(x,r)+I Q(x,r)
\]
is a smooth Fueter-regular axial function on the circularization of $U$.
There exists a slice-regular axial function $f$ on the same set such that
$\lap f=g$. It is unique up to addition of $qa+b$, $a,b\in\HH$.
\end{theorem}
\begin{proof}
Equation~\eqref{eq:axial-D}, evaluated for all $I$, gives the Vekua system
\begin{equation}\label{eq:vekua}
 P_x-Q_r-\frac{2Q}{r}=0,\qquad Q_x+P_r=0.
\end{equation}
The quaternion-valued one-form
\begin{equation}\label{eq:inverse-first}
 \omega=-\frac P2\dd x+\frac Q2\dd r
\end{equation}
is closed by the second equation. Integrating its four real components on
the simply connected set $U$ gives a function $C$ with
$C_x=-P/2$ and $C_r=Q/2$. Set $B=rC$.

Next consider
\begin{equation}\label{eq:inverse-second}
 \eta=\left(C+\frac{rQ}{2}\right)\dd x+\frac{rP}{2}\dd r.
\end{equation}
It is closed because
\[
 \partial_r\left(C+\frac{rQ}{2}\right)
 =Q+\frac{rQ_r}{2}
 =\frac{rP_x}{2}
 =\partial_x\left(\frac{rP}{2}\right).
\]
Choose a primitive $A$. Then
\[
 A_x=C+\frac{rQ}{2}=B_r,\qquad
 A_r=\frac{rP}{2}=-B_x.
\]
Thus $A+\iota B$ is a holomorphic stem on $U$; defining it on the reflected
lower set by complex conjugation gives the required stem on the symmetric
union. Its induced function $f=A+IB$ is slice regular. Formula
\eqref{eq:slice-D-lap} now gives
\[
 \lap f=\frac{2A_r}{r}
       +I\left(\frac{2B_r}{r}-\frac{2B}{r^2}\right)=P+IQ.
\]
Changing the first primitive $C$ by a constant quaternion $a$ changes $B$
by $ra$ and changes the second primitive $A$ by $xa+b$. These are precisely
the additions $qa+b$. Conversely, the kernel calculation in the proof of
Theorem~\ref{thm:fueter-map}, carried out locally on $U$, shows that the
difference of any two primitives has this form; connectedness makes the
constants global on $U$.
\end{proof}
This construction explains both the affine ambiguity and the topological
hypothesis. On a non-simply-connected planar set, closed one-forms can have
nonzero periods. The theorem does not claim a global primitive without
period conditions, nor a primitive for a general nonaxial Fueter-regular
function.

\section{What survives, what changes, and why}
\label{sec:comparison}
The preceding proofs identify two different mechanisms behind quaternionic
analogues. Slice analysis retains one-complex-variable holomorphy in each
section, but coordinates the sections through a stem and a representation
formula. Fueter analysis retains a first-order elliptic differential
operator and its fundamental solution, but no longer admits all ordinary
powers of $q$.

\begin{center}
\small
\renewcommand{\arraystretch}{1.22}
\begin{tabularx}{\textwidth}{@{}p{.19\textwidth}XX@{}}
\toprule
Classical feature & Slice-regular counterpart & Fueter-regular counterpart\\
\midrule
Cauchy formula & Planar contour, regular resolvent kernel
(\S\ref{sec:cauchy-slice}) & Three-dimensional boundary, Dirac kernel
(\S\ref{sec:fueter-cauchy})\\
Taylor theory & Powers $(q-c)^n a_n$ at real centers
(\S\ref{sec:cauchy-slice}) & Real-coordinate analyticity and derivative
estimates (\S\ref{sec:fueter-consequences})\\
Harmonicity & Biharmonic: $\lap^2f=0$
(\S\ref{sec:fueter-map}) & Harmonic: $\lap f=0$
(\S\ref{sec:fueter-consequences})\\
Zeros & Isolated points or isolated whole spheres; scalar norm controls
multiplicity (\S\ref{sec:algebra}--\ref{sec:counting}) & A nonzero function
can vanish on a plane (Example~\ref{ex:fueter-plane})\\
Maximum and Liouville & Both hold
(\S\ref{sec:maximum}) & Both hold
(\S\ref{sec:fueter-consequences})\\
Open mapping & Holds off the closure of degenerate spheres and for
symmetric open sets (\S\ref{sec:reciprocal}) & Fails in general
(Example~\ref{ex:fueter-plane})\\
Singularities & Laurent series and residues at real isolated singularities
(\S\ref{sec:laurent}) & Flux residues and a sharp $o(r^{-3})$ removability
criterion (\S\ref{sec:fueter-consequences})\\
Conformal mapping & Intrinsic slice-wise Riemann mapping, not arbitrary
four-dimensional conformality (\S\ref{sec:geometry}) & No corresponding
claim follows from the Fueter equation alone\\
\bottomrule
\end{tabularx}
\end{center}

Three cautions recur. First, a right coefficient cannot be moved to the
left without changing the problem. Second, the geometric unit for many
slice theorems is the whole conjugacy sphere, not an individual point.
Third, a quaternionic theorem must specify which regularity it uses.
With those distinctions in place, substantial parts of the classical
analytic toolkit do survive, and the Laplacian provides a concrete bridge
between the two surviving structures.

\section*{Bibliographical note}
\addcontentsline{toc}{section}{Bibliographical note}
This paper is an expository synthesis, not a claim of new priority for the
theorems. The slice-regular framework and its classical-function-theory
analogues originate in the sources cited throughout. The stem formulation,
regular algebra, and scalar norm method follow the structural viewpoint
of~\cite{GP}. The zero and open-mapping theories are treated
in~\cite{GStZeros,GStOpen}; the noncommutative Schwarz--Pick theory is
in~\cite{BS}. The real-centered singularity theory presented here is a
restricted part of the broader theory in~\cite{GPStSing}. The
Cauchy--Fueter part is derived from the divergence theorem in the spirit
of~\cite{Sudbery}, and the final elementary two-integration construction
is a quaternionic axial instance of inverse Fueter theory
\cite{CSSInverse,CPSSInverse}. The supplied proofs make these statements
independent of an appeal to unproved quaternionic machinery; the standard
scalar complex theorems listed in the opening scope statement remain
prerequisites.

\begin{thebibliography}{10}
\addcontentsline{toc}{section}{References}

\bibitem{GS}
G.~Gentili and D.~C.~Struppa,
\emph{A new theory of regular functions of a quaternionic variable},
Advances in Mathematics \textbf{216} (2007), 279--301.
\href{https://doi.org/10.1016/j.aim.2007.05.010}{doi:10.1016/j.aim.2007.05.010}.

\bibitem{GP}
R.~Ghiloni and A.~Perotti,
\emph{Slice regular functions on real alternative algebras},
Advances in Mathematics \textbf{226} (2011), 1662--1691.
\href{https://arxiv.org/abs/1008.4318}{arXiv:1008.4318}.

\bibitem{CGS}
F.~Colombo, G.~Gentili, and I.~Sabadini,
\emph{A Cauchy kernel for slice regular functions},
Annals of Global Analysis and Geometry \textbf{37} (2010), 361--378.
\href{https://doi.org/10.1007/s10455-009-9191-7}{doi:10.1007/s10455-009-9191-7};
\href{https://arxiv.org/abs/0902.4771}{arXiv:0902.4771}.

\bibitem{GStZeros}
G.~Gentili and C.~Stoppato,
\emph{Zeros of regular functions and polynomials of a quaternionic variable},
Michigan Mathematical Journal \textbf{56} (2008), 655--667.
\href{https://doi.org/10.1307/mmj/1231770366}{doi:10.1307/mmj/1231770366}.

\bibitem{GStOpen}
G.~Gentili and C.~Stoppato,
\emph{The open mapping theorem for regular quaternionic functions},
Annali della Scuola Normale Superiore di Pisa, Classe di Scienze (5)
\textbf{8} (2009), 805--815.
\href{https://doi.org/10.2422/2036-2145.2009.4.07}{doi:10.2422/2036-2145.2009.4.07};
\href{https://arxiv.org/abs/0802.3861}{arXiv:0802.3861}.

\bibitem{BS}
C.~Bisi and C.~Stoppato,
\emph{The Schwarz--Pick lemma for slice regular functions},
Indiana University Mathematics Journal \textbf{61} (2012), 297--317.
\href{https://doi.org/10.1512/iumj.2012.61.5076}{doi:10.1512/iumj.2012.61.5076};
\href{https://arxiv.org/abs/1209.2060}{arXiv:1209.2060}.

\bibitem{GPStSing}
R.~Ghiloni, A.~Perotti, and C.~Stoppato,
\emph{Singularities of slice regular functions over real alternative
$*$-algebras},
Advances in Mathematics \textbf{305} (2017), 1085--1130.
\href{https://arxiv.org/abs/1606.03609}{arXiv:1606.03609}.

\bibitem{Sudbery}
A.~Sudbery,
\emph{Quaternionic analysis},
Mathematical Proceedings of the Cambridge Philosophical Society
\textbf{85} (1979), 199--225.
\href{https://dougsweetser.github.io/Q/Stuff/pdfs/Quaternionic-analysis.pdf}{Author's paper, electronic reprint}.

\bibitem{CSSInverse}
F.~Colombo, I.~Sabadini, and F.~Sommen,
\emph{The inverse Fueter mapping theorem},
Communications on Pure and Applied Analysis \textbf{10} (2011), 1165--1181.
\href{https://doi.org/10.3934/cpaa.2011.10.1165}{doi:10.3934/cpaa.2011.10.1165}.

\bibitem{CPSSInverse}
F.~Colombo, D.~Pe\~na Pe\~na, I.~Sabadini, and F.~Sommen,
\emph{A new integral formula for the inverse Fueter mapping theorem},
Journal of Mathematical Analysis and Applications \textbf{417} (2014),
112--122.
\href{https://arxiv.org/abs/1302.0685}{arXiv:1302.0685}.

\end{thebibliography}
\end{document}
